% =====================================================================
% Fluid Dynamics and IBM Theory: nondimensionalization and similarity
% =====================================================================
\documentclass[aspectratio=169,11pt]{beamer}

% Shared style, ported from the UF Animal Sciences 2026 talk (ib-tutorials.sty at the repo root).
\makeatletter\def\input@path{{./}{../}}\makeatother
\usepackage{ib-tutorials}

\usepackage[utf8]{inputenc}
\usepackage{graphicx}
\usepackage{amsmath,amssymb,bm}
\usepackage{booktabs}
\usepackage{xcolor}
\usepackage{tcolorbox}
\tcbuselibrary{skins}
\usepackage{hyperref}

\graphicspath{{../pictures/1-0/}{pictures/1-0/}}

\newcommand{\uvec}{\bm{u}}
\newcommand{\xvec}{\bm{x}}
\newcommand{\bvec}{\bm{b}}

\title[Scaling and similarity]{Nondimensionalization and similarity}
\subtitle{Fluid, transport, and immersed-structure parameters}
\author[Miller]{Laura A. Miller}
\institute[UArizona]{Department of Mathematics, University of Arizona}
\date{Fluid Dynamics and IBM Theory \textbullet\ Module 0}

\begin{document}

% ----- 1. Title -----
% The right-hand panel of \maketutorialtitle shows the key dimensionless groups
% instead of an image: \includegraphics is redirected for this frame only.
\begin{frame}[plain]
\renewcommand{\includegraphics}[2][]{%
  \begin{punchbox}
  \centering
  \[
    \mathrm{Re}=\frac{\rho U L}{\mu},
    \qquad
    \mathrm{St}_L=\frac{fL}{U}
  \]
  \[
    \alpha=L\sqrt{\frac{\omega}{\nu}},
    \qquad
    \mathrm{Pe}=\frac{UL}{D}
  \]
  \end{punchbox}}
\maketutorialtitle[dimensionless-groups]
\end{frame}

% ----- 2. Purpose -----
\begin{frame}{Why nondimensionalize?}
\begin{columns}[T,onlytextwidth]
\begin{column}{0.54\textwidth}
Nondimensional variables let us:
\begin{itemize}\itemsep5pt
  \item compare systems with different sizes, speeds, and materials;
  \item identify which physical effects dominate;
  \item reduce the number of independent parameters;
  \item design dynamically similar experiments and simulations.
\end{itemize}
\end{column}
\begin{column}{0.42\textwidth}
\begin{punchbox}
\small
A nondimensional solution can always be converted back to physical units once the reference scales are known.
\end{punchbox}
\medskip
\soft{Choose and report the scales. A value such as $\mathrm{Re}=100$ is only reproducible when $L$ and $U$ are defined.}
\end{column}
\end{columns}
\vfill
{\scriptsize\color{DarkGray} Background: S. Vogel, \emph{Life in Moving Fluids} (1994) and D. J. Acheson, \emph{Elementary Fluid Dynamics} (1990). Definitions of the Strouhal, Womersley, and P\'eclet numbers were adapted from Wikipedia (CC BY-SA 4.0) in the original notes. Full references are at the end of the deck.\par}
\end{frame}

% ----- 3. Scales -----
\begin{frame}{Characteristic scales}
Choose scales from the physics of the problem:
\begin{center}
\renewcommand{\arraystretch}{1.18}
\begin{tabular}{lll}
\toprule
Quantity & Scale & Possible choice \\
\midrule
Length & $L$ & chord, diameter, body length \\
Velocity & $U$ & mean inflow, swim speed, tip speed \\
Time & $T=L/U$ & advective time \\
Pressure & $P_I=\rho U^2$ & inertial pressure \\
\bottomrule
\end{tabular}
\end{center}
\medskip
\begin{columns}[T,onlytextwidth]
\begin{column}{0.48\textwidth}
\begin{punchbox}
\small
Use the same $L$ and $U$ in every dimensionless group and coefficient unless a different scale is explicitly introduced.
\end{punchbox}
\end{column}
\begin{column}{0.48\textwidth}
\small
The computational domain must be large enough that moving the outer boundary farther away does not change the quantity of interest. There is no universal ``$10L$'' rule.
\end{column}
\end{columns}
\end{frame}

% ----- 4. Variables -----
\begin{frame}{Dimensionless variables and derivatives}
Define starred variables by
\[
  \xvec=L\xvec^*,
  \qquad
  \uvec=U\uvec^*,
  \qquad
  t=\frac{L}{U}t^*,
  \qquad
  p-p_0=\rho U^2 p_I^*.
\]

The differential operators transform as
\[
  \nabla=\frac{1}{L}\nabla^*,
  \qquad
  \frac{\partial}{\partial t}=\frac{U}{L}\frac{\partial}{\partial t^*},
  \qquad
  \frac{D\uvec}{Dt}=\frac{U^2}{L}
  \left[
  \frac{\partial\uvec^*}{\partial t^*}
  +(\uvec^*\!\cdot\!\nabla^*)\uvec^*
  \right].
\]

\begin{punchbox}
\small A prime can be mistaken for a derivative. Stars make the dimensionless variables visually distinct throughout the derivation.
\end{punchbox}
\end{frame}

% ----- 5. NS -----
\begin{frame}{Inertial form of incompressible Navier--Stokes}
For a Newtonian fluid with constant density $\rho$ and viscosity $\mu$,
\[
\boxed{
  \frac{\partial\uvec^*}{\partial t^*}
  +(\uvec^*\!\cdot\!\nabla^*)\uvec^*
  =-\nabla^*p_I^*
  +\frac{1}{\mathrm{Re}}\nabla^{*2}\uvec^*
  +\bvec^*
}
\]
\[
\boxed{\nabla^*\!\cdot\uvec^*=0},
\qquad
\mathrm{Re}=\frac{\rho U L}{\mu}=\frac{UL}{\nu},
\qquad
\nu=\frac{\mu}{\rho}.
\]
{\small Derivation: divide the dimensional momentum equation by $\rho U^2/L$; the viscous term becomes $\mu U/L^2\div\rho U^2/L=1/\mathrm{Re}$.}

Here $\bvec^*=(L/U^2)\bvec$ is a dimensionless body-force acceleration. If $\bvec=-g\hat{\bm\jmath}$, then $\bvec^*=-(1/\mathrm{Fr}^2)\hat{\bm\jmath}$, where $\mathrm{Fr}=U/\sqrt{gL}$ is the Froude number.
\end{frame}

% ----- 6. Re -----
\begin{frame}{Reynolds number}
\begin{columns}[T,onlytextwidth]
\begin{column}{0.52\textwidth}
\centering
\includegraphics[width=\linewidth]{Re.png}

{\scriptsize Flight speed versus $\log_{10}\mathrm{Re}$, from dust to airships. Redrawn after P.~B.~S. Lissaman, ``Low-Reynolds-number airfoils,'' \emph{Annu. Rev. Fluid Mech.} 15 (1983), 223--239, Fig.~1.}
\end{column}
\begin{column}{0.44\textwidth}
\[
  \mathrm{Re}=\frac{\text{inertial scale}}{\text{viscous scale}}
  =\frac{\rho U^2/L}{\mu U/L^2}.
\]
\begin{itemize}\itemsep5pt
  \item $\mathrm{Re}\ll1$: viscous stresses dominate; inertia is a perturbation.
  \item $\mathrm{Re}\gg1$: inertia is important; thin shear layers and separation may occur.
\end{itemize}
\smallskip
{\footnotesize Each unit of the abscissa spans a decade in Re.}
\begin{punchbox}
\small Reynolds number identifies a balance. It does not, by itself, determine whether a flow is turbulent.
\end{punchbox}
\end{column}
\end{columns}
\end{frame}

% ----- 7. Viscous scaling -----
\begin{frame}{Viscous normalization at low Reynolds number}
The inertial pressure scale is convenient when inertia is important:
\[
  \frac{D\uvec^*}{Dt^*}
  =-\nabla^*p_I^*+\frac{1}{\mathrm{Re}}\nabla^{*2}\uvec^*+\bvec^*.
\]
For creeping-flow problems, choose the viscous pressure scale
\[
  p-p_0=\frac{\mu U}{L}p_V^*.
\]
The same dimensional equation then becomes
\[
  \boxed{
  \mathrm{Re}\frac{D\uvec^*}{Dt^*}
  =-\nabla^*p_V^*+\nabla^{*2}\uvec^*+\mathrm{Re}\,\bvec^*
  }.
\]
\begin{punchbox}
\small Both forms are valid at every Reynolds number. They satisfy $p_V^*=\mathrm{Re}\,p_I^*$, and the inertially scaled body acceleration acquires the same factor; there is no mathematical switch at $\mathrm{Re}=1$.
\end{punchbox}
\end{frame}

% ----- 8. Time and frequency -----
\begin{frame}{Dimensionless time and frequency}
The advective time scale $L/U$ gives
\[
  t^*=\frac{tU}{L},
  \qquad
  f^*=\frac{fL}{U},
  \qquad
  \omega^*=\frac{\omega L}{U}=2\pi f^*.
\]
\begin{itemize}\itemsep6pt
  \item One dimensionless cycle has period $1/f^*$.
  \item Matching $f^*$ preserves the ratio of oscillation time to advection time.
  \item Pumping, flapping, and pulsatile-flow comparisons must use the same definitions of $L$ and $U$.
\end{itemize}
\begin{punchbox}
\small A dimensional frequency cannot be compared across geometrically or dynamically different systems until it is combined with the relevant time scale.
\end{punchbox}
\end{frame}

\begin{frame}{Frequency scaling and the Strouhal number}
For vortex shedding or another oscillation with frequency $f$,
\[
  \boxed{\mathrm{St}_L=\frac{fL}{U}}.
\]
This compares the oscillation time scale $1/f$ with the advective time scale $L/U$.

For an oscillating wake, wing, or propulsor,
\[
  \boxed{\mathrm{St}_A=\frac{fA}{U}},
\]
where $A$ is the stated excursion scale. If $A=L$, then $\mathrm{St}_A=\mathrm{St}_L$.

\begin{punchbox}
\small State whether the characteristic length is a chord, diameter, wake width, or motion amplitude. Different choices define different Strouhal numbers.
\end{punchbox}
\soft{\scriptsize The number is named for V. Strouhal, who studied vortex shedding from wires in 1878. Definition adapted from Wikipedia, ``Strouhal number'' (CC BY-SA 4.0), as in the original notes.}
\end{frame}

\begin{frame}{Strouhal number in swimming and flying}
For a swimming or flying animal, a common kinematic definition is
\[
  \boxed{\mathrm{St}_A=\frac{fA}{U}},
\]
where $f$ is stroke or tail-beat frequency, $A$ is a stated excursion amplitude, and $U$ is mean forward speed.
\begin{itemize}\itemsep6pt
  \item $\mathrm{St}_A$ organizes vortex-growth and shedding regimes for pitching and heaving propulsors.
  \item Many cruising swimmers and flyers in the Taylor et al. data lie in $0.2<\mathrm{St}_A<0.4$.
  \item The interval is empirical, not a universal efficiency optimum for every body, gait, or amplitude convention.
\end{itemize}
\soft{\scriptsize State whether $A$ is peak or peak-to-peak excursion; Taylor et al.\ use the peak-to-peak excursion of the wingtip or tail tip. Taylor, Nudds, and Thomas, \emph{Nature} 425 (2003), 707--711.}
\end{frame}

% ----- 9. Womersley -----
\begin{frame}{Womersley number}
For oscillatory viscous flow with angular frequency $\omega$,
\[
  \boxed{\alpha=L\sqrt{\frac{\omega}{\nu}}}.
\]
\begin{columns}[T,onlytextwidth]
\begin{column}{0.50\textwidth}
\begin{itemize}\itemsep5pt
  \item $\alpha\ll1$: viscosity communicates the oscillation across the length scale during each cycle.
  \item $\alpha\gg1$: unsteady inertia dominates away from a thin oscillatory boundary layer.
\end{itemize}
\end{column}
\begin{column}{0.46\textwidth}
\[
  \delta_\omega\sim\sqrt{\frac{2\nu}{\omega}},
  \qquad
  \alpha=\sqrt{2}\,\frac{L}{\delta_\omega}.
\]
If $\mathrm{Re}$ and $\mathrm{St}_L=fL/U$ use the same $L$ and $U$,
\[
  \boxed{\alpha^2=2\pi\,\mathrm{Re}\,\mathrm{St}_L}.
\]
\end{column}
\end{columns}
\soft{\footnotesize If one group uses pipe radius and another uses diameter, the relation acquires the corresponding factor.}

\soft{\scriptsize Interpretation: $\alpha^2=\omega L^2/\nu$ is the ratio of transient inertia, $\rho\omega U$, to viscous force, $\mu U/L^2$. Named for J. R. Womersley (\emph{J. Physiol.} 127, 1955). Definition adapted from Wikipedia, ``Womersley number'' (CC BY-SA 4.0), as in the original notes.}
\end{frame}

\begin{frame}{Illustrative Womersley numbers in canine vessels}
The original tutorial used the following values for a canine cardiovascular system at a heart rate near $2\,\mathrm{Hz}$. Here the length scale in $\alpha=L\sqrt{\omega/\nu}$ is the vessel radius, $L=R$.
\begin{columns}[T,onlytextwidth]
\begin{column}{0.49\textwidth}
\begin{center}
\renewcommand{\arraystretch}{1.16}
\begin{tabular}{lr}
\toprule
Vessel & $\alpha$ \\
\midrule
Ascending aorta & $13.2$ \\
Descending aorta & $11.5$ \\
Abdominal aorta & $8.0$ \\
Femoral artery & $3.5$ \\
Carotid artery & $4.4$ \\
\bottomrule
\end{tabular}
\end{center}
\end{column}
\begin{column}{0.49\textwidth}
\begin{center}
\renewcommand{\arraystretch}{1.16}
\begin{tabular}{lr}
\toprule
Vessel & $\alpha$ \\
\midrule
Arterioles & $0.04$ \\
Capillaries & $0.005$ \\
Venules & $0.035$ \\
Inferior vena cava & $8.8$ \\
Main pulmonary artery & $15$ \\
\bottomrule
\end{tabular}
\end{center}
\end{column}
\end{columns}
\begin{punchbox}
\small These values illustrate the strong vessel-size dependence of $\alpha$. They are not universal physiological constants: radius, frequency, and viscosity must accompany any reported value.
\end{punchbox}
\soft{\scriptsize Values after W. W. Nichols and M. F. O'Rourke, \emph{McDonald's Blood Flow in Arteries}, 5th ed. (Hodder Arnold, 2005). Further sources on the next slide.}
\end{frame}

\begin{frame}{Canine Womersley numbers: a worked check and sources}
\textbf{Order-of-magnitude check} for the ascending aorta. Take an illustrative radius $R\approx0.75\,\mathrm{cm}$, heart rate $f=2\,\mathrm{Hz}$, and blood kinematic viscosity $\nu\approx4\times10^{-6}\,\mathrm{m^2\,s^{-1}}$:
\[
  \omega=2\pi f\approx12.6\,\mathrm{s^{-1}},
  \qquad
  \alpha=R\sqrt{\frac{\omega}{\nu}}
  \approx7.5\times10^{-3}\,\mathrm{m}\times\sqrt{\frac{12.6}{4\times10^{-6}}}\,\mathrm{m^{-1}}
  \approx13,
\]
consistent with the tabulated $13.2$. With the diameter in place of the radius, the same vessel would give $\alpha\approx27$; always state which length is used.

\medskip
\begin{punchbox}
\small Because $\alpha\propto R$, a factor of $\sim10^3$ in radius between the aorta and a capillary produces the factor of $\sim10^3$ in $\alpha$ seen in the table.
\end{punchbox}
\soft{\scriptsize Sources: the same ten canine values are commonly cited from the table of typical canine vascular parameters in C. G. Caro, T. J. Pedley, R. C. Schroter, and W. A. Seed, \emph{The Mechanics of the Circulation} (Oxford Univ. Press, 1978), as reproduced in Wikipedia, ``Womersley number'' (CC BY-SA 4.0); table number to be confirmed. The original tutorial's Womersley material was adapted from that Wikipedia article.}
\end{frame}

% ----- 10. Peclet -----
\begin{frame}{Péclet number}
For a scalar concentration $c$ with molecular diffusivity $D$,
\[
  \frac{\partial c}{\partial t}+\uvec\cdot\nabla c=D\nabla^2c.
\]
With $c-c_0=Cc^*$ and the same $L$, $U$, and $L/U$ scales,
\[
  \frac{\partial c^*}{\partial t^*}
  +\uvec^*\cdot\nabla^*c^*
  =\frac{1}{\mathrm{Pe}}\nabla^{*2}c^*,
  \qquad
  \boxed{\mathrm{Pe}=\frac{UL}{D}}.
\]
\begin{columns}[T,onlytextwidth]
\begin{column}{0.48\textwidth}
\hi{$\mathrm{Pe}\ll1$}: diffusion dominates transport.
\end{column}
\begin{column}{0.48\textwidth}
\hi{$\mathrm{Pe}\gg1$}: advection dominates transport.
\end{column}
\end{columns}
\medskip
For mass transfer, $\mathrm{Pe}=\mathrm{Re}\,\mathrm{Sc}$; for heat transfer, $\mathrm{Pe}=\mathrm{Re}\,\mathrm{Pr}$, provided all groups use the same $L$ and $U$.

\soft{\scriptsize Here $\mathrm{Sc}=\nu/D$ and $\mathrm{Pr}=\nu/\kappa$, with $\kappa$ the thermal diffusivity. Named for J. C. E. P\'eclet. Definitions adapted from Wikipedia, ``P\'eclet number'' (CC BY-SA 4.0), as in the original notes.}
\end{frame}

% ----- 11. Force scales -----
\begin{frame}{Force scales for a 2-D model}
Suppose $\mathcal F$ denotes a force resultant per unit out-of-plane depth, with units $\mathrm{N\,m^{-1}}$. For example, it may be the integral of boundary traction over a length of order $L$.
\begin{center}
\renewcommand{\arraystretch}{1.3}
\begin{tabular}{lll}
\toprule
Normalization & Stress scale & Resultant scale \\
\midrule
Inertial & $\rho U^2$ & $\mathcal F_I=\rho U^2L$ \\
Viscous & $\mu U/L$ & $\mathcal F_V=\mu U$ \\
\bottomrule
\end{tabular}
\end{center}
\[
  \mathcal F_I^*=\frac{\mathcal F}{\rho U^2L},
  \qquad
  \mathcal F_V^*=\frac{\mathcal F}{\mu U},
  \qquad
  \boxed{\mathcal F_V^*=\mathrm{Re}\,\mathcal F_I^*}.
\]
\begin{punchbox}
\small A local force per unit boundary length and per unit depth has units $\mathrm{N\,m^{-2}}$ and scales with $\rho U^2$ or $\mu U/L$. Begin with the force law and its measure.
\end{punchbox}
\end{frame}

% ----- 12. Structural coefficients -----
\begin{frame}{Dimensionless bending coefficient in two dimensions}
Under a two-dimensional per-unit-depth convention, let $k_b$ have units
\[
  [k_b]=\mathrm{N\,m}.
\]
The inertial and viscous normalizations are
\[
  \boxed{k_b^{(I)}=\frac{k_b}{\rho U^2L^3}},
  \qquad
  \boxed{k_b^{(V)}=\frac{k_b}{\mu UL^2}}.
\]
They describe the same dimensional coefficient and satisfy
\[
  k_b^{(V)}=\mathrm{Re}\,k_b^{(I)}.
\]
\begin{punchbox}
\small There is no mathematical switch at $\mathrm{Re}=1$. Choose the normalization that exposes the balance of interest, then use it consistently.
\end{punchbox}
\end{frame}

\begin{frame}{Dimensionless stretching coefficient}
Let $k_s$ denote a force-like fiber coefficient per unit out-of-plane depth:
\[
  [k_s]=\mathrm{N\,m^{-1}}.
\]
The corresponding nondimensional coefficients are
\[
  \boxed{k_s^{(I)}=\frac{k_s}{\rho U^2L}},
  \qquad
  \boxed{k_s^{(V)}=\frac{k_s}{\mu U}},
\]
with
\[
  k_s^{(V)}=\mathrm{Re}\,k_s^{(I)}.
\]
\begin{punchbox}
\small A point-spring stiffness and three-dimensional line, surface, or volume coefficients have different units. Begin with the exact force law and integration measure.
\end{punchbox}
\end{frame}

\begin{frame}{Viscosity and Reynolds-number matching}
Because
\[
  \mathrm{Re}=\frac{\rho UL}{\mu}=\frac{UL}{\nu},
\]
one can change Reynolds number by changing $U$, $L$, $\rho$, or $\mu$.
\begin{itemize}\itemsep6pt
  \item Varying $\mu$ is common in simulations and laboratory scale models.
  \item Changing viscosity may also change other material or transport parameters.
  \item Reporting only $\mu$ does not establish dynamic similarity; report $\mathrm{Re}$ and the chosen $L$ and $U$.
\end{itemize}
\begin{punchbox}
\small When an oscillation or scalar field matters, matching Reynolds number alone is insufficient. Match the relevant Strouhal, Womersley, Péclet, and structural groups as well.
\end{punchbox}
\end{frame}

% ----- 13. EI -----
\begin{frame}{Flexural rigidity and effective stiffness}
\begin{columns}[T,onlytextwidth]
\begin{column}{0.50\textwidth}
For a physical three-dimensional beam,
\[
  \boxed{B=EI},
\]
where
\begin{itemize}\itemsep4pt
  \item $E$ is Young's modulus $[\mathrm{N\,m^{-2}}]$;
  \item $I$ is a second moment of area $[\mathrm{m^4}]$;
  \item $EI$ is flexural rigidity $[\mathrm{N\,m^2}]$.
\end{itemize}
\end{column}
\begin{column}{0.46\textwidth}
\begin{punchbox}
\small
Flexural rigidity is a material-and-section property. The force/displacement stiffness also depends on beam length, loading, and boundary conditions; commonly $K_{\mathrm{eff}}\sim EI/L_b^3$.
\end{punchbox}
\medskip
A 2-D model may use rigidity per unit out-of-plane depth, with units $\mathrm{N\,m}$. State that convention explicitly.
\end{column}
\end{columns}
\end{frame}

\begin{frame}{Approximate Young's modulus: compliant materials}
\footnotesize
\setlength{\tabcolsep}{3pt}
\renewcommand{\arraystretch}{1.0}
\begin{center}
\begin{tabular}{@{}p{0.55\textwidth}rr@{}}
\toprule
Material & GPa & psi \\
\midrule
Rubber, small strain & $0.01$--$0.1$ & $1{,}450$--$14{,}504$ \\
Low-density polyethylene & $0.11$--$0.45$ & $16{,}000$--$65{,}000$ \\
Diatom frustules, largely hydrated silica & $0.35$--$2.77$ & $50{,}000$--$400{,}000$ \\
PTFE & $0.5$ & $72{,}500$ \\
High-density polyethylene & $0.8$ & $116{,}000$ \\
Bacteriophage capsids & $1$--$3$ & $145{,}000$--$435{,}000$ \\
Polypropylene & $1.5$--$2$ & $218{,}000$--$290{,}000$ \\
Polyethylene terephthalate & $2.2$--$2.7$ & $320{,}000$--$390{,}000$ \\
Nylon & $2$--$4$ & $290{,}000$--$580{,}000$ \\
Polystyrene & $3$--$3.5$ & $440{,}000$--$510{,}000$ \\
Medium-density fiberboard & $4$ & $580{,}000$ \\
Oak, along the grain & $11$ & $1.60\times10^6$ \\
Human cortical bone & $14$ & $2.03\times10^6$ \\
\bottomrule
\end{tabular}
\end{center}
\begin{punchbox}
\small These are order-of-magnitude reference values. Modulus depends on composition, orientation, hydration, strain, and test protocol.
\end{punchbox}
{\tiny\color{DarkGray} Conversion: $1\,\mathrm{GPa}\approx145{,}038\,\mathrm{psi}$. Values follow the table ``Approximate Young's modulus for various materials'' in Wikipedia, ``Young's modulus'' (CC BY-SA 4.0), which cites a primary source for each entry; source to be confirmed.\par}
\end{frame}

\begin{frame}{Approximate Young's modulus: stiffer materials}
\footnotesize
\setlength{\tabcolsep}{3pt}
\renewcommand{\arraystretch}{1.08}
\begin{center}
\begin{tabular}{@{}p{0.55\textwidth}rr@{}}
\toprule
Material & GPa & psi \\
\midrule
Glass-reinforced polyester matrix & $17.2$ & $2.49\times10^6$ \\
Aromatic peptide nanotubes & $19$--$27$ & $2.76$--$3.92\times10^6$ \\
High-strength concrete & $30$ & $4.35\times10^6$ \\
Carbon-fiber-reinforced plastic, 50/50 biaxial & $30$--$50$ & $4.35$--$7.25\times10^6$ \\
Hemp fiber & $35$ & $5.08\times10^6$ \\
Magnesium & $45$ & $6.53\times10^6$ \\
Glass & $50$--$90$ & $7.25$--$13.1\times10^6$ \\
Flax fiber & $58$ & $8.41\times10^6$ \\
Aluminum & $69$ & $10.0\times10^6$ \\
Mother-of-pearl, largely calcium carbonate & $70$ & $10.2\times10^6$ \\
Aramid & $70.5$--$112.4$ & $10.2$--$16.3\times10^6$ \\
Tooth enamel, largely calcium phosphate & $83$ & $12.0\times10^6$ \\
\bottomrule
\end{tabular}
\end{center}
\begin{punchbox}
\small Use a source appropriate to the actual material and loading direction before converting $E$ into flexural rigidity $EI$.
\end{punchbox}
{\tiny\color{DarkGray} Conversion: $1\,\mathrm{GPa}\approx145{,}038\,\mathrm{psi}$. Values follow the same Wikipedia ``Young's modulus'' table (CC BY-SA 4.0); source to be confirmed.\par}
\end{frame}

% ----- 14. I -----
\begin{frame}{Second moment of area}
For a cross-sectional area $A$ and an in-plane reference axis $\ell$,
\[
  \boxed{I_\ell=\iint_A r_\perp^2\,dA},
\]
where $r_\perp$ is the perpendicular distance from the axis to the area element $dA$.

\begin{columns}[T,onlytextwidth]
\begin{column}{0.44\textwidth}
For the Cartesian axes,
\[
  I_x=\iint_A y^2\,dA,
  \qquad
  I_y=\iint_A x^2\,dA.
\]
\end{column}
\begin{column}{0.52\textwidth}
The polar second moment about the perpendicular $z$ axis is
\[
  J_z=\iint_A(x^2+y^2)\,dA=I_x+I_y.
\]
\end{column}
\end{columns}
\begin{punchbox}
\small $I_x$ and $I_y$ control bending about the corresponding axes. $J_z$ is a polar area moment and should not be confused with a mass moment of inertia.
\end{punchbox}
\soft{\scriptsize The axis is written $\ell$ to avoid confusion with the flexural rigidity $B=EI$. For an axis parallel to a centroidal axis at distance $d$, the parallel-axis theorem gives $I_\ell=I_{\mathrm{centroid}}+Ad^2$.}
\end{frame}

\begin{frame}{Rectangle with centroid at the origin}
\begin{columns}[T,onlytextwidth]
\begin{column}{0.54\textwidth}
For width $b$ in the $x$ direction and height $h$ in the $y$ direction,
\[
  I_x
  =\int_{-b/2}^{b/2}\int_{-h/2}^{h/2}y^2\,dy\,dx
  =\frac{bh^3}{12},
\]
\[
  I_y
  =\int_{-h/2}^{h/2}\int_{-b/2}^{b/2}x^2\,dx\,dy
  =\frac{b^3h}{12},
\]
\[
  J_z=I_x+I_y
  =\frac{bh}{12}(b^2+h^2).
\]
\end{column}
\begin{column}{0.42\textwidth}
\centering
\includegraphics[width=\linewidth]{IJ.png}

{\scriptsize Rectangular cross-section of width $b$ and height $h$ with centroid $C$ (schematic). Schematic: instructor's course notes.}

\medskip
{\small Bending about $x$ (load along $y$) uses $I_x\propto h^3$: a plank is far stiffer on edge than flat.}
\end{column}
\end{columns}
\end{frame}

% ----- 15. Conversion -----
\begin{frame}{Converting structural parameters}
\begin{enumerate}\itemsep7pt
  \item Select and document one set of $L$ and $U$.
  \item Choose inertial or viscous normalization for the full coupled model.
  \item Convert the desired dimensionless coefficient back to physical units.
\end{enumerate}

For example, with inertial normalization,
\[
  k_{\mathrm{spring}}
  =k_{\mathrm{spring},I}^*\,\rho U^2L,
  \qquad
  k_{\mathrm{beam}}
  =k_{\mathrm{beam},I}^*\,\rho U^2L^3.
\]
{\small In the notation of the earlier slides, $k_{\mathrm{spring},I}^*=k_s^{(I)}$ and $k_{\mathrm{beam},I}^*=k_b^{(I)}$ (2-D, per unit depth).}

\begin{punchbox}
\small A coefficient does not determine deformation by itself. Geometry, loading, boundary conditions, and the constitutive law all matter, so report a sensitivity study rather than a universal ``good'' value.
\end{punchbox}
\end{frame}

\begin{frame}{Beam-stiffness starting estimate}
\small
With inertial normalization,
\[
  k_b=k_b^{(I)}\rho U^2L^3.
\]
The original course heuristic used $k_b^{(I)}\approx10$ as a trial value that produced modest bending, roughly a few degrees, in a particular model.

\begin{itemize}\itemsep2pt
  \item Use the same characteristic velocity and length as the fluid scaling; a peak velocity may be appropriate for a peak-load estimate.
  \item In practice: try $U_{\max}$ for the velocity scale and the boundary length for $L$; these are rough estimates and should be reported with the choice made. If $\mathrm{Re}$ was defined with a different velocity $U$ (same $L$), the two coefficients for one dimensional $k_b$ satisfy $k_b^{(I)}\big|_{U_{\max}}=k_b^{(I)}\big|_{U}\,(U/U_{\max})^2$.
  \item Compare the resulting $k_b$ with a physical rigidity such as $EI$ under the model's per-depth convention.
  \item Sweep $k_b^{(I)}$ and demonstrate convergence of the deformation or performance metric.
\end{itemize}
\begin{punchbox}
\small The numerical value $10$ is a model-specific starting guess, not a universal relation between dimensionless stiffness and bending angle.
\end{punchbox}
\end{frame}

% ----- 16. Loads -----
\begin{frame}{Physical load estimates}
At moderate or high Reynolds number, estimate a total hydrodynamic load from a physical area:
\[
  F_{\mathrm{char}}\sim\frac12 C_D\rho U^2A.
\]
For a 2-D model per unit depth, use the corresponding projected length instead of area. At low Reynolds number, use a geometry-specific viscous drag law.

\begin{columns}[T,onlytextwidth]
\begin{column}{0.50\textwidth}
For a target displacement $\delta_{\max}$,
\[
  K_{\mathrm{eff}}\sim\frac{F_{\mathrm{char}}}{\delta_{\max}}
\]
provides an order-of-magnitude starting point.
\end{column}
\begin{column}{0.46\textwidth}
\begin{punchbox}
\small
Do not use the Eulerian grid spacing $\Delta x$ as a physical area or mass unless you are explicitly constructing a discrete per-node quantity. Otherwise the physical model changes when the grid is refined.
\end{punchbox}
\end{column}
\end{columns}
\end{frame}

\begin{frame}{Spring-stiffness starting estimate}
If $F_b$ is the characteristic load supported by a spring-like constraint and $\delta_{\max}$ is the largest acceptable displacement,
\[
  \boxed{K_{\mathrm{eff}}\sim\frac{F_b}{\delta_{\max}}}.
\]
Choose $\delta_{\max}$ from the physical tolerance, such as maximum displacement from a target point or resting length.

\begin{columns}[T,onlytextwidth]
\begin{column}{0.49\textwidth}
At moderate or high Reynolds number,
\[
  F_b\sim\tfrac12 C_D\rho U^2A.
\]
\end{column}
\begin{column}{0.47\textwidth}
At low Reynolds number, use a geometry-specific viscous drag law rather than a universal grid-based estimate.
\end{column}
\end{columns}
\begin{punchbox}
\small Convert the resulting stiffness using the units of the implemented force law. A node, line, surface, and volume spring require different coefficients.
\end{punchbox}
\end{frame}

% ----- 17. Damping -----
\begin{frame}{Damping in a coupled fluid--structure system}
For one linear degree of freedom,
\[
  m\ddot{x}+c\dot{x}+kx=0,
  \qquad
  c_{\mathrm{crit}}=2\sqrt{km}.
\]
This formula is a useful analogy, but an immersed structure has distributed modes and frequency-dependent fluid loading.

\begin{itemize}\itemsep5pt
  \item A massless boundary still experiences added mass and hydrodynamic damping.
  \item The effective mass is generally not $\rho(\Delta x)^d$.
  \item Select penalty damping with a modal estimate or solver-specific stability and convergence tests.
\end{itemize}

\begin{punchbox}
\small Report the damping law and demonstrate that the physical response and time-step restriction are adequately resolved.
\end{punchbox}
\soft{\scriptsize Critical damping, $c=c_{\mathrm{crit}}$, returns the oscillator to rest fastest without overshoot; the damping ratio is $\zeta_d=c/c_{\mathrm{crit}}$. Background: Wikipedia, ``Damping'' (CC BY-SA 4.0), linked in the original notes.}
\end{frame}

% ----- 18. Efficiency -----
\begin{frame}{Efficiency at steady swimming or flying}
\small
In a statistically steady or time-periodic state, mean kinetic energy does not grow, although it may oscillate instantaneously. Input power still balances drag and energizes the wake.

{\small A thermodynamic efficiency (useful work, i.e.\ change in kinetic or potential energy, divided by energy input) would therefore be zero for every steady swimmer. Efficiency must instead be defined through a chosen performance metric, as below.}

\begin{columns}[T,onlytextwidth]
\begin{column}{0.48\textwidth}
\begin{block}{Constrained propulsor}
A common definition is
\[
  \eta_{\mathrm{prop}}=\frac{TU}{P_{\mathrm{in}}},
\]
where $T$ is mean thrust.
\end{block}
\end{column}
\begin{column}{0.48\textwidth}
\begin{block}{Self-propelled body}
Mean net force is zero. One alternative is
\[
  \eta_{\mathrm{QP}}=\frac{D_{\mathrm{tow}}U}{P_{\mathrm{in}}}.
\]
Here $D_{\mathrm{tow}}$ is the drag from a matched towing test.
\end{block}
\end{column}
\end{columns}

\begin{punchbox}
\small State the system boundary, useful-power definition, and treatment of negative actuator work; different efficiency measures answer different questions.
\end{punchbox}
\soft{\scriptsize For a review of thrust, drag, and efficiency definitions in swimming, see Smits, \emph{J. Fluid Mech.} 874 (2019), P1.}
\end{frame}

% ----- 19. COT -----
\begin{frame}{Cost of transport}
\small
If $E_{\mathrm{in}}$ is the input energy used to travel distance $d$,
\[
  C=\frac{E_{\mathrm{in}}}{d}
  =\frac{\overline{P}_{\mathrm{in}}}{\overline{U}}
  \quad [\mathrm{J\,m^{-1}}=\mathrm{N}].
\]

Common normalizations are
\begin{columns}[T,onlytextwidth]
\begin{column}{0.41\textwidth}
\[
  C_m=\frac{E_{\mathrm{in}}}{m d}
  =\frac{\overline{P}_{\mathrm{in}}}{m\overline{U}}
  \quad [\mathrm{J\,kg^{-1}\,m^{-1}}],
\]
\end{column}
\begin{column}{0.57\textwidth}
\[
  C^*=\frac{\overline{P}_{\mathrm{in}}}{m g\overline{U}}
  \quad \text{(dimensionless, weight-normalized)}.
\]
\end{column}
\end{columns}

For a distributed actuator, compute input power from the work rate using a stated sign convention, for example
\[
  P_{\mathrm{in}}(t)=\int_{\Gamma}\bm{f}_{\mathrm{act}}\cdot\uvec_b\,d\Gamma,
\]
where $\bm{f}_{\mathrm{act}}$ is the actuator force density on the structure $\Gamma$ and $\uvec_b$ is the boundary velocity. Average $P_{\mathrm{in}}$ over a whole number of cycles before forming $C$.

\soft{\scriptsize Recall: work $=$ force $\times$ distance [N\,m]; power $=$ force $\times$ velocity [N\,m\,s$^{-1}$]. The mass-specific cost $C_m$ is the form compared across animals by K. Schmidt-Nielsen, \emph{Science} 177 (1972), 222--228; the weight-normalized $C^*$ is the ``specific resistance'' of G. Gabrielli and T. von K\'arm\'an, \emph{Mech. Eng.} 72 (1950), 775--781.}
\end{frame}

% ----- 20. Summary -----
\begin{frame}{Nondimensionalization checklist}
\begin{enumerate}\itemsep8pt
  \item Define $L$, $U$, and every additional material or transport scale.
  \item Substitute the scaled variables into the dimensional equations.
  \item Choose one normalization that exposes the balance of interest.
  \item Check every dimensionless coefficient by units and limiting cases.
  \item Preserve the same definitions when comparing cases or redimensionalizing results.
\end{enumerate}
\begin{punchbox}
\centering
\small Dynamic similarity requires matching every dimensionless group that materially affects the problem, not only Reynolds number.
\end{punchbox}
\end{frame}

% ----- 21. References -----
\begin{frame}{References}
\small
\begin{itemize}\itemsep4pt
  \item D. J. Acheson, \emph{Elementary Fluid Dynamics}. Clarendon Press, 1990.
  \item S. Vogel, \emph{Life in Moving Fluids}, 2nd ed. Princeton University Press, 1994.
  \item G. K. Taylor, R. L. Nudds, and A. L. R. Thomas, ``Flying and swimming animals cruise at a Strouhal number tuned for high power efficiency,'' \emph{Nature} 425 (2003), 707--711, \href{https://doi.org/10.1038/nature02000}{doi:10.1038/nature02000}.
  \item J. R. Womersley, ``Method for the calculation of velocity, rate of flow and viscous drag in arteries when the pressure gradient is known,'' \emph{J. Physiol.} 127 (1955), 553--563, \href{https://doi.org/10.1113/jphysiol.1955.sp005276}{doi:10.1113/jphysiol.1955.sp005276}.
  \item A. J. Smits, ``Undulatory and oscillatory swimming,'' \emph{J. Fluid Mech.} 874 (2019), P1, \href{https://doi.org/10.1017/jfm.2019.284}{doi:10.1017/jfm.2019.284}.
  \item C. S. Peskin, ``The immersed boundary method,'' \emph{Acta Numerica} 11 (2002), \href{https://doi.org/10.1017/S0962492902000077}{doi:10.1017/S0962492902000077}.
  \item B. E. Griffith and X. Y. Luo, ``Hybrid finite difference/finite element version of the immersed boundary method,'' \emph{Int. J. Numer. Meth. Biomed. Engng.} 33 (2017), \href{https://doi.org/10.1002/cnm.2888}{doi:10.1002/cnm.2888}.
\end{itemize}
\end{frame}

\begin{frame}{References (continued)}
\footnotesize
\begin{itemize}\itemsep4pt
  \item P. B. S. Lissaman, ``Low-Reynolds-number airfoils,'' \emph{Annu. Rev. Fluid Mech.} 15 (1983), 223--239,
  \href{https://doi.org/10.1146/annurev.fl.15.010183.001255}{doi:10.1146/annurev.fl.15.010183.001255}. Fig.~1 redrawn on the Reynolds-number slide.
  \item W. W. Nichols and M. F. O'Rourke, \emph{McDonald's Blood Flow in Arteries}, 5th ed. Hodder Arnold, 2005.
  \item C. G. Caro, T. J. Pedley, R. C. Schroter, and W. A. Seed, \emph{The Mechanics of the Circulation}. Oxford University Press, 1978 (2nd ed., Cambridge University Press, 2012). Canine Womersley values.
  \item K. Schmidt-Nielsen, ``Locomotion: energy cost of swimming, flying, and running,'' \emph{Science} 177 (1972), 222--228,
  \href{https://doi.org/10.1126/science.177.4045.222}{doi:10.1126/science.177.4045.222}.
  \item G. Gabrielli and T. von K\'arm\'an, ``What price speed? Specific power required for propulsion of vehicles,'' \emph{Mechanical Engineering} 72 (1950), 775--781.
  \item Wikipedia articles ``Strouhal number,'' ``Womersley number,'' ``P\'eclet number,'' ``Damping,'' and ``Young's modulus'' (CC BY-SA 4.0), \url{https://en.wikipedia.org}; definitions and the modulus table were adapted from these in the original notes.
\end{itemize}
\begin{punchbox}
\small Figure credits: \texttt{Re.png} redrawn after Lissaman (1983), Fig.~1; \texttt{IJ.png} schematic from the instructor's course notes.
\end{punchbox}
\end{frame}

\end{document}
